\documentclass[12pt]{article}
\usepackage[T1]{fontenc}
\usepackage[utf8]{inputenc}
\usepackage[polish]{babel}
\usepackage{lmodern}
\usepackage[a4paper,margin=2.4cm]{geometry}
\usepackage{xcolor}
\usepackage{graphicx}
\usepackage{amsmath,amssymb}
\usepackage{booktabs}
\usepackage{array}
\usepackage{enumitem}
\usepackage{hyperref}
\usepackage{tcolorbox}
\usepackage{fancyhdr}
\definecolor{WEIIPrimary}{RGB}{0,130,150}
\definecolor{WEIIPrimaryDark}{RGB}{0,86,100}
\definecolor{WEIIPrimaryLight}{RGB}{229,246,248}
\definecolor{PRzNavy}{RGB}{31,63,142}
\definecolor{RLTextDark}{RGB}{34,34,34}
\definecolor{RLGray}{RGB}{90,90,90}
\definecolor{RLLightGray}{RGB}{245,246,248}
\definecolor{RLWarning}{RGB}{184,86,0}
\definecolor{RLExample}{RGB}{52,112,64}
\definecolor{RLCodeBg}{RGB}{248,248,248}

\usepackage{listings}
\usepackage{xcolor}

\lstdefinestyle{rlpython}{%
    language=Python,
    basicstyle=\ttfamily\scriptsize,
    keywordstyle=\color{WEIIPrimaryDark}\bfseries,
    commentstyle=\color{RLGray}\itshape,
    stringstyle=\color{RLExample},
    numbers=left,
    numberstyle=\tiny\color{RLGray},
    numbersep=6pt,
    showstringspaces=false,
    breaklines=true,
    tabsize=4,
    frame=single,
    rulecolor=\color{WEIIPrimary!35},
    backgroundcolor=\color{RLCodeBg},
    captionpos=b
}

\lstset{style=rlpython}


\hypersetup{colorlinks=true,linkcolor=WEIIPrimaryDark,urlcolor=WEIIPrimaryDark,citecolor=WEIIPrimaryDark}
\setlength{\headheight}{15pt}
\pagestyle{fancy}
\fancyhf{}
\lhead{Uczenie się ze wzmocnieniem}
\rhead{WEiI PRz}
\cfoot{\thepage}

\newtcolorbox{taskbox}[1]{colback=WEIIPrimaryLight,colframe=WEIIPrimary,title=\textbf{#1}}
\newtcolorbox{notebox}[1]{colback=RLLightGray,colframe=RLGray,title=\textbf{#1}}
\newtcolorbox{warninghandout}[1]{colback=RLWarning!8,colframe=RLWarning,title=\textbf{#1}}

\newcommand{\handouttitle}[3]{%
\begin{center}
{\Large\bfseries #1}\\[0.2cm]
{\large #2}\\[0.15cm]
{\small #3}
\end{center}
\vspace{0.4cm}
\hrule
\vspace{0.6cm}
}

\newcommand{\EnableLabNumbering}{%
  \renewcommand{\thesection}{L\arabic{section}}%
  \renewcommand{\thesubsection}{\thesection.\arabic{subsection}}%
  \renewcommand{\thesubsubsection}{\thesubsection.\arabic{subsubsection}}%
}


\title{Raport projektowy\\Uczenie się ze wzmocnieniem}
\author{Patryk Przybyło}
\date{2026/2027}

\begin{document}
\maketitle
\tableofcontents
\newpage

\section{Informacje podstawowe}

\begin{center}
\begin{tabular}{p{0.34\textwidth}p{0.56\textwidth}}
\toprule
Element & Informacja \\
\midrule
Imię i nazwisko & Patryk Przybyło \\
Numer albumu & 179597 \\
Kierunek/specjalność & EF-DI/AI \\
Tytuł projektu & Symulacja ramienia sortującego \\
Link do repozytorium & \url{https://github.com/xcxcp/ucznie_ze_wzmocnieniem} \\
Data oddania & Do uzupełnienia \\
\bottomrule
\end{tabular}
\end{center}

\section{Temat i cel projektu}

\subsection{Opis problemu}

Problemem jest sterowanie manipulatorem o wielu stopniach swobody w środowisku symulacyjnym silnika fizycznego.
Robot jest odzwzorowany na podstawie rzeczywistego ramienia a nauczona przez agenta strategia w symulacji ma być przeniesiona na rzeczywiste sterowanie.

\subsection{Cel agenta}

Na podstawie pozycji obiektu agent ma tak ustawić ramię, aby chwytak znajdował się nad obiektem w taki sposób, aby możliwe było chwycenie obiektu.
Przy czym ma on unikać zderzeń z innymi obiektami i podłożem oraz nie wykonywać ruchów naruszających fizyczne ograniczenia manipulatora.

\subsection{Uzasadnienie zastosowania uczenia się ze wzmocnieniem}

Sterowania ramieniem robota  można rozpatrywać jako problem uczenia się ze wzmocnieniem, ponieważ robot podejmuje serię powiązanych ze sobą decyzji, które wpływają na jego dalsze działanie oraz na stan otoczenia. 
W trakcie realizacji zadania robot musi kolejno podejmować decyzje dotyczące między innymi kierunku i zakresu ruchu poszczególnych przegubów, położenia ramienia oraz wysunięcia siłownika liniowego.
Decyzje te tworzą sekwencję prowadzącą od początkowego położenia robota do wykonania zadania.
Każde działanie agenta zmienia stan środowiska. Przykładowo obrót podstawy, zmiana kąta ramienia lub wysunięcie siłownika wpływają na aktualne położenie chwytaka.
W sterowaniu robotem istotne jest osiągnięcie końcowego celu, a nie tylko uzyskanie korzyści z pojedynczego ruchu. Przykładowo szybkie zbliżenie chwytaka do obiektu nie oznacza jeszcze sukcesu, jeżeli chwytak nie będzie ustawiony w odpowiedniej pozycji.
Dlatego agent powinien uwzględniać konsekwencje swoich działań w kolejnych etapach zadania.
Zastosowanie agenta uczącego się jest uzasadnione ze względu na złożoność sterowania robotem posiadającym kilka stopni swobody oraz konieczność dostosowywania ruchów do zmiennego położenia.
Ręczne opracowanie reguł sterowania dla wszystkich możliwych sytuacji może być trudne i mało elastyczne. Agent uczący się ze wzmocnieniem może natomiast zdobywać doświadczenie poprzez wielokrotne wykonywanie prób w środowisku symulacyjnym.

\section{Formalizacja problemu RL/MDP}

\subsection{Agent i środowisko}

Agentem jest algorytm uczenia się ze wzmocnieniem odpowiedzialny za sterowanie ramieniem robota. 
Jego zadaniem jest wybieranie kolejnych ruchów prowadzących do znalezienia się chywtaka nad obiektem.
Środowiskiem jest model robota oraz otoczenia przygotowany w symulatorze MuJoCo. Obejmuje ono:
\begin{itemize}
    \item obrotową podstawę robota,
    \item człony ramienia i jego przeguby,
    \item siłownik liniowy,
    \item nadgarstek z niezależnymi osiami obrotu,
    \item chwytak wyposażony w trzy zsynchronizowane palce,
    \item sorotowane obiekty,
    \item plansza.
\end{itemize}
Agent i środowisko współpracują w następującej pętli:
\begin{enumerate}
    \item Środowisko udostępnia agentowi aktualną obserwację.
    \item Agent analizuje obserwację i wybiera akcję.
    \item Akcja jest przekazywana do środowiska symulacyjnego.
    \item Symulator wykonuje zadany ruch, uwzględniając mechanikę robota i oddziaływania z otoczeniem.
    \item Środowisko zwraca nową obserwację, wartość nagrody oraz informację, czy epizod został zakończony.
    \item Agent wykorzystuje otrzymane informacje do aktualizacji swojej strategii podczas treningu.
\end{enumerate}

\subsection{Stan lub obserwacja}

Agent ma dostęp do informacji o kątach obrotu podstawy, przegubów ramienia i nadgarstk, wysunięciu siłownika liniowego oraz położeniu obiektów.
Nie ma natomiast informacji bezpośredniej o położeniu poszczególnych części robota ani gdzie doszło do kolizcji.

\begin{center}
\begin{tabular}{p{0.30\textwidth}p{0.60\textwidth}}
\toprule
Składowa stanu/obserwacji & Znaczenie \\
\midrule
$a_b$ & Kąt obrotu ramienia \\
$a_r$ & Kąt między ramieniem a podstawą \\
$a_l$ & Wysunięcie siłownika liniowego \\
$a_e$ & Kąt zgięcia łokcia\\
$a_n$ & Kąt obrotu nadgarskta wokół własnej osi\\
$a_c$ & Kąt obrotu chwytaka względem nadgastak\\
$x_i$, $y_i$, $z_i$ & Połóżenie obiektu\\
\bottomrule
\end{tabular}
\end{center}

\subsection{Akcje}

Zbiór akcji określa, jakie działania agent może wykonywać w środowisku. W projektowanym robocie akcje dotyczą sterowania położeniem poszczególnych osi mechanizmu oraz chwytakiem.
Przyjęto, że agent będzie sterował wartościami zadanymi dla napędów, a nie bezpośrednio sygnałami elektrycznymi lub PWM. 
Takie rozwiązanie pozwala oddzielić proces podejmowania decyzji przez agenta od niskopoziomowego sterowania silnikami.
W proponowanym rozwiązaniu przestrzeń akcji jest ciągła, agent wybiera wartości zadane ruchu w określonych zakresach.
Wartości akcji przekraczające dopuszczalne zakresy są odrzucane przez warstwę sterowania. 
Dzięki temu agent nie może wymusić ruchu prowadzącego do mechanicznego uszkodzenia robota.

\begin{center}
\begin{tabular}{p{0.30\textwidth}p{0.60\textwidth}}
\toprule
Akcja & Opis \\
\midrule
$a_b$ & Kąt obrotu podstawy \\
$a_e$ & Kąt obortu łokcia \\
$a_n$ & Kąt obrotu nadgarstka \\
$a_c$ & Kąt obrotu chwytaka\\
$a_l$ & Wysunięcie siłownika liniowego\\
\bottomrule
\end{tabular}
\end{center}

\subsection{Funkcja nagrody}

Funkcja nagrody określa, które zachowania agenta są pożądane. 
Jej zadaniem jest zachęcanie do poprawnego ustawienia chwytaka, a jednocześnie ograniczanie niepotrzebnych ruchów i niebezpiecznych zachowań.

\begin{equation}
    r_t = w_u R_{\text{ustawienie}} + w_p(d_{t-1}-d_t) + w_k P_{\text{kolizcja}} + w_n P_{\text{niedozwolony ruch}} + w_t
\end{equation}

Gdzie:
\begin{itemize}
    \item $d_t$ - odległość chwytaka od aktualnego celu
    \item $R_{\text{ustawienie}}$ - nagroda za poprawne ustawienie chywtaka
    \item $P_{\text{kolizcja}}$ - kara za kolizje
    \item $P_{\text{niedozwolony ruch}}$ - kara za niedozwolony ruch
    \item $w_t$ - kara za krok
    \item $w_u$, $w_p$, $w_k$, $w_n$ - współczynniki określające znaczenie poszczególnych składników
\end{itemize}

\begin{center}
\begin{tabular}{p{0.30\textwidth}p{0.60\textwidth}}
\toprule
Składnik & Proponowana wartość \\
\midrule
$w_u$ & 100\\
$w_p$ & 1 \\
$w_k$ & -50 \\
$w_n$ & -50\\
$w_t$ & -0.01\\
\bottomrule
\end{tabular}
\end{center}

Nagroda za poprawne ustawienie jest największa, ponieważ stanowi główny cel zadania.
Kara za kolizję i niedozwolony ruch ma ograniczać zachowania niebezpieczne.
Ważne jest, aby funkcja nagrody nie zachęcała agenta do osiągania krótkoterminowych korzyści kosztem całego zadania. 
Przykładowo samo zbliżenie chwytaka do obiektu nie może być traktowane jako sukces.


\subsection{Warunki zakończenia epizodu}

Epizod oznacza pojedynczą próbę wykonania zadania. Rozpoczyna się od ustawienia robota i obiektu w określonych pozycjach początkowych, a kończy po wystąpieniu jednego z określonych zdarzeń.
Epizod zostaje zakończony w następujących sytuacjach:
\begin{itemize}
    \item Sukces: agent poprawnie umieścił chwytak nad obiektem
    \item Przekroczenie limitu kroków: agent nie wykonał zadania w określonej liczbie kroków.
\end{itemize}

\subsection{Założenia i ograniczenia środowiska}

Na potrzeby projektu przyjęto uproszczony model środowiska, który umożliwia przeprowadzenie treningu agenta w symulatorze MuJoCo.
Najważniejsze założenia i uproszczenia:
\begin{itemize}
    \item Rrobot jest odwzorowany w MuJoCo jako układ sztywnych elementów połączonych przegubami obrotowymi i przesuwnymi.
    \item Serwomechanizmy są traktowane jako napędy realizujące zadane pozycje. Nie odwzorowuje się w pełni wszystkich zjawisk występujących w rzeczywistych serwach, takich jak luzy przekładni, opóźnienia komunikacji czy zmiany momentu w zależności od warunków pracy.
    \item Agent dostaje dokładne pozycje obiektów, gdzie w rzeczywistym robocie pozycje są wyznaczane z kamer
    \item Model symulacyjny nie odwzorowuje idealnie rzeczywistej konstrukcji. Mogą występować różnice w masie, bezwładności, tarciu, sztywności elementów
\end{itemize}

Przyjęte uproszczenia pozwalają jednak rozpocząć trening w kontrolowanym środowisku, stopniowo zwiększać jego złożoność i przygotować algorytm do późniejszego wykorzystania w rzeczywistym systemie sortowania.

\section{Implementacja}

\subsection{Struktura repozytorium}

Należy krótko opisać strukturę repozytorium i najważniejsze pliki projektu.

\begin{verbatim}
repozytorium/
|-- src/
|-- experiments/
|-- report/
|-- results/
`-- README.md
\end{verbatim}

\subsection{Środowisko}

Należy opisać sposób implementacji środowiska, w szczególności funkcje inicjalizacji epizodu oraz wykonania kroku środowiska.

\subsection{Agent bazowy}

Należy opisać agenta losowego, heurystycznego albo inny baseline wykorzystany do porównania z agentem uczącym się.

\subsection{Agent uczący się}

Należy opisać zastosowany algorytm RL lub DRL, jego parametry oraz sposób implementacji. W przypadku Q-tab należy opisać reprezentację tablicy wartości. W przypadku DRL należy opisać wejścia modelu, wyjścia modelu, architekturę oraz sposób uczenia.

\section{Plan eksperymentów}

Należy opisać eksperymenty zaplanowane w celu oceny działania agenta.

\begin{center}
\begin{tabular}{p{0.24\textwidth}p{0.30\textwidth}p{0.30\textwidth}}
\toprule
Eksperyment & Co jest zmieniane? & Co jest mierzone? \\
\midrule
E1 & Do uzupełnienia & Do uzupełnienia \\
E2 & Do uzupełnienia & Do uzupełnienia \\
E3 & Do uzupełnienia & Do uzupełnienia \\
\bottomrule
\end{tabular}
\end{center}

W tej sekcji należy wskazać:
\begin{itemize}
    \item liczbę epizodów uczenia,
    \item liczbę powtórzeń eksperymentu albo uzasadnienie braku powtórzeń,
    \item parametry uczenia,
    \item metryki oceny,
    \item warianty środowiska, map lub funkcji nagrody,
    \item sposób porównania z agentem bazowym.
\end{itemize}

\section{Wyniki}

W tej sekcji należy przedstawić wyniki eksperymentów. Wyniki powinny być pokazane w postaci tabel, wykresów lub czytelnych zestawień.

\begin{center}
\begin{tabular}{p{0.28\textwidth}p{0.20\textwidth}p{0.20\textwidth}p{0.20\textwidth}}
\toprule
Wariant & Średni zwrot & Skuteczność & Długość epizodu \\
\midrule
Agent bazowy & Do uzupełnienia & Do uzupełnienia & Do uzupełnienia \\
Agent RL/DRL & Do uzupełnienia & Do uzupełnienia & Do uzupełnienia \\
\bottomrule
\end{tabular}
\end{center}

\begin{figure}[h]
    \centering
    \fbox{\parbox{0.82\textwidth}{\centering Miejsce na wykres procesu uczenia}}
    \caption{Przykładowy wykres procesu uczenia agenta.}
\end{figure}

\section{Analiza i wnioski}

Należy zinterpretować wyniki. W tej części należy odpowiedzieć m.in. na pytania:
\begin{itemize}
    \item Czy agent uczący się działa lepiej niż agent bazowy?
    \item Jak zmieniało się zachowanie agenta w trakcie uczenia?
    \item Które elementy środowiska były najtrudniejsze?
    \item Czy funkcja nagrody prowadziła do oczekiwanego zachowania?
    \item Czy wystąpiły przypadki niepożądanego zachowania agenta?
    \item Jakie są najważniejsze ograniczenia rozwiązania?
\end{itemize}

\section{Instrukcja uruchomienia}

Należy opisać sposób uruchomienia projektu od czystego repozytorium. Przykładowo:

\begin{verbatim}
git clone https://github.com/...
cd nazwa-repozytorium
python -m venv .venv
.venv\Scripts\activate
pip install -r requirements.txt
python src/main.py
\end{verbatim}

Należy również wskazać, jak odtworzyć eksperymenty oraz gdzie zapisywane są wyniki.

\section{Wsparcie AI}

W tej sekcji należy obowiązkowo opisać wykorzystanie narzędzi sztucznej inteligencji. Jeżeli student nie korzystał z takich narzędzi, należy wpisać jednoznaczną informację, że nie korzystano ze wsparcia AI.

\begin{center}
\begin{tabular}{p{0.13\textwidth}p{0.19\textwidth}p{0.21\textwidth}p{0.27\textwidth}}
\toprule
Data & Narzędzie/model & Cel użycia & Transkrypt / link / zakres \\
\midrule
Do uzupełnienia & Do uzupełnienia & Do uzupełnienia & Do uzupełnienia \\
Do uzupełnienia & Do uzupełnienia & Do uzupełnienia & Do uzupełnienia \\
\bottomrule
\end{tabular}
\end{center}

Należy zapisać transkrypty rozmów z modelami AI wraz z datą oraz nazwą modelu. Jeżeli rozmowa odbywała się poza narzędziami programistycznymi, należy dodatkowo udostępnić całą rozmowę poprzez link, o ile użyte narzędzie umożliwia takie udostępnienie.

Student powinien także wyjaśnić, które elementy wygenerowane lub zasugerowane przez AI zostały samodzielnie sprawdzone, zmodyfikowane albo odrzucone.

\section{Literatura}

W tej sekcji należy wskazać wykorzystane źródła: książki, artykuły naukowe, dokumentację bibliotek, materiały wykładowe oraz inne źródła użyte podczas pracy nad projektem.

\begin{thebibliography}{9}
\bibitem{sutton_barto}
R. S. Sutton, A. G. Barto, \textit{Reinforcement Learning: An Introduction}, 2nd ed., MIT Press, 2018.
\end{thebibliography}

\end{document}
